% =============================================================================
% Spanish content definitions — content/es.tex
% =============================================================================
% Defines all user-visible strings for the Spanish resume variant.
% Every macro name must match content/en.tex exactly.
% =============================================================================

% Document language (consumed by \usepackage[\documentlanguage]{babel})
\newcommand{\documentlanguage}{spanish}

% =============================================================================
% Section headings
% =============================================================================
\newcommand{\secProfessionalSummary}{Resumen Profesional}
\newcommand{\secTechnicalSkills}{Habilidades Técnicas}
\newcommand{\secProfessionalExperience}{Experiencia Profesional}
\newcommand{\secKeyProjects}{Proyectos}
\newcommand{\secEducationalTools}{Herramientas Educativas y Demostraciones}
\newcommand{\secProfessionalSkills}{Habilidades Profesionales}
\newcommand{\secConferences}{Presentaciones en Conferencias}
\newcommand{\secEducation}{Educación}

% =============================================================================
% Subsection titles
% =============================================================================
\newcommand{\subNTRU}{Protocolo de Intercambio de Llaves Post-Cuántico [C++, NTRU]}
\newcommand{\subAES}{Biblioteca de Cifrado AES [C/C++, CMake, GTest, Estándares NIST]}
\newcommand{\subPixelLab}{Pixel Lab - Kit para Criptoanálisis de Imágenes [Python, NumPy, Bash Scripting]}
\newcommand{\subThesis}{Protocolo de Comunicación Segura [Protocolos Criptográficos, Arquitectura de Sistemas, Seguridad Post-Cuántica]}
\newcommand{\subEduTools}{Sitio Web de Portafolio de Seguridad [HTML, CSS, JavaScript, Criptografía Web]}

% =============================================================================
% Project date lines
% =============================================================================
\newcommand{\dateNTRU}{Marzo 2023 -- Junio 2025 (En mantenimiento y actualización activa)}
\newcommand{\dateAES}{Enero 2024 -- Junio 2025 (En mantenimiento y actualización activa)}
\newcommand{\datePixelLab}{Enero 2025 -- Presente (Desarrollo activo)}
\newcommand{\dateThesis}{Agosto 2023 -- Julio 2025 (Investigación de Tesis de Maestría)}
\newcommand{\dateEduTools}{Julio 2025 -- Presente (Mantenimiento activo)}

% =============================================================================
% Header labels
% =============================================================================
\newcommand{\labelLocation}{Ubicación}
\newcommand{\labelContact}{Contacto}
\newcommand{\labelPortfolio}{Portafolio}
\newcommand{\labelGitHub}{GitHub}
\newcommand{\labelLinkedIn}{LinkedIn}

% =============================================================================
% Role title subtitles (header and footer)
% =============================================================================
\newcommand{\roleTitleCrypto}{Ingeniero en Criptografía}
\newcommand{\roleTitleSecurity}{Ingeniero de Seguridad (Esp. en Criptografía)}
\newcommand{\roleTitleSoftware}{Ingeniero de Software (Esp. en Seguridad)}
\newcommand{\roleTitleApplied}{Criptógrafo Aplicado}

% =============================================================================
% Professional Summary — body paragraphs
% =============================================================================
\newcommand{\summaryCrypto}{%
	Especialista en criptografía con amplia experiencia en sistemas criptográficos post-cuánticos e implementaciones de nivel productivo. Experto en criptografía reticular NTRU, cifrado AES y estándares de seguridad conformes a NIST. Historial comprobado en el diseño de protocolos resistentes a ataques cuánticos, optimización de primitivas criptográficas y validación rigurosa de seguridad. Sólida formación matemática respaldada por investigación de maestría en protocolos de comunicación segura. En búsqueda de roles desafiantes en ingeniería criptográfica para desarrollar infraestructura de seguridad de próxima generación e impulsar la adopción de criptografía post-cuántica.%
}

\newcommand{\summarySecurity}{%
	Ingeniero de seguridad especializado en implementaciones criptográficas y seguridad post-cuántica, con experiencia práctica en el desarrollo de sistemas de cifrado de nivel productivo. Experiencia comprobada implementando bibliotecas de cifrado AES conformes a NIST, desarrollando protocolos de intercambio de claves resistentes a ataques cuánticos (NTRU) y optimizando el rendimiento criptográfico. Sólida formación en desarrollo de software seguro, pruebas de seguridad, criptoanálisis y validación de cumplimiento normativo. Desarrollo de herramientas web educativas y mantenimiento de portafolio público que demuestra conceptos criptográficos. En búsqueda de roles en ingeniería de seguridad para diseñar infraestructura resiliente, implementar protocolos criptográficos y fortalecer la postura de seguridad organizacional.%
}

\newcommand{\summarySoftware}{%
	Ingeniero de software con especialización en seguridad y criptografía, con experiencia en el desarrollo de sistemas de alto rendimiento con bases en seguridad. Habilidades en programación de sistemas en C/C++, arquitectura de software, marcos de pruebas automatizadas y optimización de rendimiento. Desarrolló una biblioteca de cifrado AES de nivel productivo y un criptosistema post-cuántico NTRU con pruebas exhaustivas y documentación completa. Fuerte enfoque en código limpio, mantenibilidad y buenas prácticas de ingeniería. En búsqueda de roles en ingeniería de software para construir sistemas seguros y escalables donde el conocimiento criptográfico represente una ventaja competitiva.%
}

\newcommand{\summaryApplied}{%
	Criptógrafo aplicado que conecta la criptografía teórica con la implementación práctica, con investigación de maestría en protocolos de comunicación resistentes a ataques cuánticos y software criptográfico de calidad productiva. Experto en criptografía reticular (NTRU), cifrado simétrico (AES) y diseño de protocolos criptográficos. Investigación publicada en criptografía post-cuántica con sólidas bases matemáticas. Desarrolla herramientas criptográficas de código abierto y recursos educativos que demuestran conceptos de seguridad complejos. En búsqueda de roles en criptografía aplicada en laboratorios de investigación, startups enfocadas en seguridad o equipos de I\&D que impulsen la teoría criptográfica hacia soluciones desplegables.%
}

% =============================================================================
% Technical Skills — bullet content
% =============================================================================
\newcommand{\skillsCryptoSecShared}{\textbf{Criptografía y Seguridad:} Criptografía Post-Cuántica (NTRU, Reticular), Mecanismos de Encapsulamiento de Llaves (KEM), RSA, Criptografía de Curva Elíptica (ECC), AES, Modos de Cifrado por Bloques (CBC, CTR, ECB, OFB), Firmas Digitales, Intercambio de Llaves, Criptoanálisis, Funciones Hash (SHA-2, SHA-3).}

\newcommand{\skillsProgrammingLangSoftware}{\textbf{Lenguajes de Programación:} C/C++ (Programación de Sistemas, Optimización), Python (Scripting, Automatización), Java, Bash Scripting, JavaScript (Desarrollo Web), HTML/CSS.}

\newcommand{\skillsSecurityEngShared}{\textbf{Ingeniería de Seguridad:} Diseño de Protocolos Criptográficos, Pruebas y Validación de Seguridad, Cumplimiento de Estándares NIST (FIPS 197, SP 800-38A, FIPS 180-4), Modelado de Amenazas (STRIDE), Métricas de Seguridad (Entropía, Chi-Cuadrada, Análisis de Correlación).}

\newcommand{\skillsSoftwareEngSoftware}{\textbf{Ingeniería de Software:} Arquitectura de Sistemas, Optimización y Perfilado de Rendimiento, Marcos de Pruebas Automatizadas (Unitarias, Integración, Sistema), Pipelines CI/CD.}

\newcommand{\skillsResearchApplied}{\textbf{Investigación y Análisis:} Diseño y Análisis de Protocolos Criptográficos, Pruebas de Seguridad y Verificación Formal, Modelado Matemático, Análisis de Complejidad Algorítmica, Escritura Académica y Publicaciones, Presentaciones en Conferencias, Desarrollo de Contenido Educativo.}

\newcommand{\skillsProgrammingLangCryptoSec}{\textbf{Lenguajes de Programación:} C/C++, Python, Java, Bash Scripting, JavaScript, HTML/CSS.}

\newcommand{\skillsDevToolsSoftware}{\textbf{Herramientas y Prácticas de Desarrollo:} Git (Control de Versiones), Make/CMake/Vite (Sistemas de Construcción), Docker (Contenedorización), GDB/Valgrind (Análisis de Memoria), , GTest/Vitest (Pruebas de Unidad, Integración y Sistema).}

\newcommand{\skillsSecurityToolsCryptoSec}{\textbf{Herramientas y Bibliotecas de Seguridad:} OpenSSL, Biblioteca GNU de Aritmética de Múltiple Precisión (GMP), Vectores de Prueba y Suites de Validación NIST, Herramientas de Análisis Estadístico (Entropía, Chi-Cuadrada, Correlación), Módulo de Seguridad por Hardware (HSM).}

\newcommand{\skillsCryptoLibsApplied}{\textbf{Bibliotecas y Herramientas Criptográficas:} OpenSSL, GMP (GNU Multiple Precision), Validación de Algoritmos Criptográficos NIST, Suites de Pruebas Estadísticas, LaTeX (Documentación Matemática), Herramientas de Visualización Criptográfica.}

\newcommand{\skillsSystemsShared}{\textbf{Sistemas y Plataformas:} Linux (Primaria), Windows, Interfaz de Línea de Comandos (CLI).}

\newcommand{\skillsDocCollabSoftware}{\textbf{Documentación y Colaboración:} Documentación Técnica, Documentación de API, Buenas Prácticas de Documentación de Código, Colaboración Interfuncional, Transferencia de Conocimiento.}

\newcommand{\skillsDocStdsCryptoSec}{\textbf{Documentación y Estándares:} Escritura Técnica, Documentación de Estándares NIST, Especificaciones de Protocolos de Seguridad, LaTeX, Presentaciones en Conferencias, Desarrollo de Material Educativo.}

\newcommand{\skillsAcademicApplied}{\textbf{Comunicación Académica:} LaTeX (Composición Tipográfica Profesional), Publicaciones Académicas, Presentaciones en Conferencias, Seminarios Técnicos (más de 21 presentaciones), Creación de Contenido Educativo, Documentación de Investigación.}

\newcommand{\skillsLanguagesShared}{\textbf{Idiomas:} Español (Nativo), Inglés (Dominio profesional completo), Alemán (Básico).}

% =============================================================================
% Professional Experience — title line
% =============================================================================
\newcommand{\expTitle}{\textbf{Asistente de Investigación (Criptografía y Seguridad)} $|$ Instituto Politécnico Nacional, Ciudad de México $|$ Enero 2022 -- Julio 2025}

% Opening statements
\newcommand{\expOpeningCrypto}{Lideró el desarrollo de software criptográfico para investigación en criptografía post-cuántica, implementando protocolos conformes a NIST y realizando análisis de seguridad rigurosos. \textbf{Logró una mejora de rendimiento de 20x en la generación de claves y operaciones de cifrado NTRU}.}

\newcommand{\expOpeningSecurity}{Lideró el desarrollo de software criptográfico para investigación en criptografía post-cuántica, implementando protocolos de seguridad de calidad productiva. \textbf{Entregó operaciones criptográficas 20x más rápidas}, permitiendo el despliegue práctico de cifrado resistente a ataques cuánticos en entornos con recursos limitados.}

\newcommand{\expOpeningSoftware}{Lideró el desarrollo de software criptográfico para investigación en criptografía post-cuántica, implementando código en C/C++ con marcos de pruebas y documentación. \textbf{Logró una mejora de rendimiento de 20x en operaciones criptográficas} mediante mejoras algorítmicas.}

\newcommand{\expOpeningApplied}{Realizó investigación en criptografía aplicada sobre sistemas post-cuánticos, conectando el diseño criptográfico teórico con implementaciones prácticas de software. \textbf{Redujo la complejidad de la multiplicación de polinomios de $O(n^4)$ a $O(n^2)$} mediante la explotación de la estructura del anillo, demostrando la aplicabilidad práctica de la criptografía reticular.}

% AES optimization bullets
\newcommand{\expAesCrypto}{Optimizó la biblioteca de cifrado AES \textbf{con almacenamiento de claves HSM basado en PKCS\#11}, logrando \textbf{un incremento de 2x en el rendimiento sobre hardware estándar} mediante optimización de bajo nivel.}

\newcommand{\expAesSecurity}{Optimizó la implementación de cifrado AES logrando \textbf{una mejora de rendimiento de 2x, aumentando significativamente el throughput en operaciones de datos seguros a gran escala}.}

\newcommand{\expAesSoftware}{Diseñó y optimizó la biblioteca de cifrado AES \textbf{con almacenamiento de llaves HSM basado en PKCS\#11}, manteniendo conformidad total con NIST FIPS 197.}

\newcommand{\expAesApplied}{Optimizó la implementación de cifrado AES logrando \textbf{una reducción de latencia del 50\% mediante patrones de acceso a memoria conscientes de caché y análisis algorítmico}, demostrando la traducción efectiva de principios de optimización teóricos a sistemas criptográficos productivos.}

% SHA shared bullet
\newcommand{\expShaShared}{Implementó el algoritmo hash SHA-512 con soporte para multiplicación de números de 128 bits, \textbf{logrando conformidad total con NIST FIPS 180-4}.}

% Validation bullets
\newcommand{\expValidationCrypto}{Desarrolló un marco de validación criptográfica integral que incorpora vectores de prueba NIST y métricas estadísticas de seguridad. \textbf{Validó correctitud al 100\%} cubriendo múltiples modos de cifrado (ECB, CBC, CTR, OFB).}

\newcommand{\expValidationSecurity}{Desarrolló suites de pruebas de seguridad incorporando vectores de prueba estándar NIST (FIPS 197, SP 800-38A). \textbf{Logró cobertura de pruebas al 100\% con validación automatizada de propiedades criptográficas}.}

\newcommand{\expValidationSoftware}{Construyó marcos de pruebas automatizadas y pipelines CI/CD para la validación de protocolos criptográficos. \textbf{Implementó más de 200 casos de prueba automatizados}.}

\newcommand{\expValidationApplied}{Diseñó metodología de validación combinando vectores de prueba NIST con análisis estadístico (entropía, chi-cuadrada, correlación). \textbf{Validó la calidad de aleatoriedad criptográfica en más de 10,000 operaciones de cifrado}, estableciendo confianza empírica en las propiedades de seguridad de la implementación.}

% Analysis bullets
\newcommand{\expAnalysisCrypto}{Realizó análisis de seguridad incluyendo evaluación de entropía, pruebas chi-cuadrada y análisis de correlación. \textbf{Verificó aleatoriedad criptográfica con un nivel de confianza del 99.9\%}.}

\newcommand{\expAnalysisSecurity}{Realizó análisis de seguridad incluyendo evaluación de entropía y pruebas estadísticas para verificar la aleatoriedad criptográfica. \textbf{Identificó y resolvió 3 vulnerabilidades de implementación}.}

\newcommand{\expAnalysisSoftware}{Diseñó una base de código modular y reutilizable. \textbf{Redujo el tiempo de integración en un 60\%} mediante interfaces bien definidas y documentación en línea extensa.}

\newcommand{\expAnalysisApplied}{Publicó resultados de investigación a través de presentaciones académicas y materiales educativos. \textbf{Impartió más de 21 seminarios técnicos alcanzando a más de 500 estudiantes e investigadores}, traduciendo eficazmente la teoría criptográfica compleja en contenido educativo accesible.}

% Shared bullets
\newcommand{\expEduMaterialsShared}{Produjo materiales educativos audiovisuales adecuados para presentaciones en conferencias y clases universitarias. \textbf{Generó más de 15 videos educativos y demostraciones interactivas} comunicando conceptos criptográficos avanzados a audiencias técnicas diversas.}

\newcommand{\expSeminarsShared}{Presentó resultados de investigación sobre criptografía post-cuántica en \textbf{más de 21 seminarios técnicos}, comunicando conceptos complejos de seguridad a audiencias de pregrado y posgrado.}

\newcommand{\expResearchProductShared}{\textbf{Producto de investigación: artículo científico sobre criptografía post-cuántica e implementaciones eficientes} (actualmente en revisión editorial previo a publicación).}

% =============================================================================
% NTRU project bullets
% =============================================================================
\newcommand{\ntruOpeningCrypto}{Implementó un Mecanismo de Encapsulamiento de Llaves (KEM) resistente a ataques cuánticos utilizando criptografía reticular NTRU, proporcionando seguridad frente a ataques de computadoras cuánticas.}

\newcommand{\ntruOpeningSecurity}{Implementó un Mecanismo de Encapsulamiento de Llaves (KEM) resistente a ataques cuánticos utilizando criptografía reticular NTRU, diseñado para resistir tanto ataques clásicos como cuánticos.}

\newcommand{\ntruOpeningSoftware}{Implementó un Mecanismo de Encapsulamiento de Llaves (KEM) post-cuántico utilizando NTRU.}

\newcommand{\ntruOpeningApplied}{Implementó un Mecanismo de Encapsulamiento de Llaves post-cuántico basado en NTRU reticular, traduciendo la criptografía reticular teórica en software validado. \textbf{Demostró la viabilidad práctica de la criptografía post-cuántica con generación de claves en 2ms y cifrado en menos de un milisegundo en hardware de uso general}.}

\newcommand{\ntruOptimizationCrypto}{Optimizó operaciones polinomiales utilizando propiedades algebraicas de la estructura del anillo. \textbf{Redujo la complejidad computacional de $O(n^4)$ a $O(n^2)$ manteniendo garantías de seguridad demostrables}, logrando una mejora de rendimiento de 20x en la generación de claves (de 100ms a 5ms).}

\newcommand{\ntruOptimizationSecurity}{Optimizó operaciones polinomiales utilizando propiedades de la estructura del anillo. \textbf{Logró una mejora de rendimiento de 20x reduciendo el tiempo de generación de claves de 100ms a 5ms}, haciendo la criptografía post-cuántica práctica para aplicaciones en tiempo real y entornos con recursos limitados.}

\newcommand{\ntruOptimizationSoftware}{Diseñó algoritmos de operaciones polinomiales aprovechando estructura del anillo. \textbf{Redujo la complejidad computacional de $O(n^4)$ a $O(n^2)$}, mejorando el throughput en un 2000\%.}

\newcommand{\ntruOptimizationApplied}{Optimizó la multiplicación polinomial en anillos de polinomios truncados mediante la explotación de la estructura algebraica. \textbf{Redujo la complejidad algorítmica de $O(n^4)$ a $O(n^2)$ con análisis formal de complejidad}, demostrando la aplicación efectiva del álgebra computacional a la optimización criptográfica.}

\newcommand{\ntruValidationCrypto}{Desarrolló un marco de validación de seguridad que incluye \textbf{análisis de entropía (logrando $H > 7.99$ bits/byte), pruebas chi-cuadrada ($\chi^2$ dentro de intervalos de confianza del 95\%) y análisis de correlación}.}

\newcommand{\ntruValidationSecurity}{Desarrolló un marco integral de validación de seguridad que incluye análisis de entropía, pruebas estadísticas y análisis de correlación, garantizando la correctitud y seguridad de la implementación.}

\newcommand{\ntruValidationSoftware}{Construyó una arquitectura modular con subsistemas dedicados para perfilado de rendimiento, pruebas de seguridad y utilidades de depuración.}

\newcommand{\ntruValidationApplied}{Diseñó una metodología de validación estadística que combina análisis teórico de seguridad con pruebas empíricas. \textbf{Estableció un protocolo de validación riguroso verificado mediante revisión por pares}, proporcionando un marco para validar implementaciones criptográficas basadas en retículos.}

% =============================================================================
% AES project bullets
% =============================================================================
\newcommand{\aesOpeningCrypto}{Diseñó e implementó una biblioteca de cifrado AES \textbf{con almacenamiento de claves HSM basado en PKCS\#11} que soporta múltiples modos de cifrado por bloques (CBC, CTR, ECB, OFB). \textbf{Logró conformidad total con NIST verificada mediante vectores de prueba oficiales}.}

\newcommand{\aesOpeningSecurity}{Diseñó y desarrolló una biblioteca de cifrado AES conforme a NIST FIPS 197 \textbf{con almacenamiento de claves HSM basado en PKCS\#11} que soporta múltiples modos de cifrado por bloques para protección de datos versátil. \textbf{Validó conformidad al 100\% con estándares NIST mediante cobertura de pruebas exhaustiva}.}

\newcommand{\aesOpeningSoftware}{Desarrolló una biblioteca de cifrado AES con conformidad a NIST FIPS 197 y arquitectura modular.}

\newcommand{\aesOpeningApplied}{Implementó una biblioteca de cifrado AES conforme a NIST FIPS 197. \textbf{Logró validación teórica y empírica mediante conformidad con la especificación formal}.}

\newcommand{\aesComplianceCrypto}{Implementó un marco de validación criptográfica \textbf{logrando conformidad al 100\% con los vectores de prueba NIST de los estándares FIPS 197 y SP 800-38A en todos los tamaños de clave (128/192/256 bits)}.}

\newcommand{\aesComplianceSecurity}{Implementó un marco de pruebas que valida conformidad al 100\% con los vectores de prueba NIST de los estándares FIPS 197 y SP 800-38A. \textbf{Logró cero vulnerabilidades de seguridad mediante validación sistemática}.}

\newcommand{\aesComplianceSoftware}{Implementó un marco de pruebas integral \textbf{con conformidad al 100\% con los vectores de prueba NIST} y suites de pruebas automatizadas.}

\newcommand{\aesComplianceApplied}{Validó la correctitud de la implementación mediante comparación sistemática con los vectores de prueba estándar NIST (FIPS 197, SP 800-38A). \textbf{Demostró alineación entre las especificaciones teóricas y la implementación en más de 300 casos de validación}, estableciendo confianza empírica en la correctitud criptográfica.}

\newcommand{\aesMetricsCrypto}{Construyó un módulo de análisis de métricas de calidad de cifrado que mide la aleatoriedad mediante cálculo de entropía, pruebas de distribución estadística y análisis de correlación.}

\newcommand{\aesMetricsSecurity}{Construyó un módulo de métricas que mide la calidad del cifrado mediante análisis de aleatoriedad y validación estadística.}

\newcommand{\aesMetricsSoftware}{Desarrolló un manejador de archivos compatible con imágenes, archivos de texto y binarios. \textbf{Procesó archivos de hasta 100MB sin fugas de memoria verificadas mediante Valgrind}.}

\newcommand{\aesMetricsApplied}{Implementó herramientas de análisis estadístico para la evaluación de calidad del cifrado, aplicando teoría de la información y teoría de probabilidad. \textbf{Estableció un marco de métricas cuantitativas utilizado en investigación revisada por pares}, validando propiedades prácticas de aleatoriedad criptográfica mediante medición empírica.}

% =============================================================================
% Pixel Lab project bullets
% =============================================================================
\newcommand{\pixellabOpeningCrypto}{Desarrolló un kit de herramientas en Python para validación de aleatoriedad criptográfica mediante criptoanálisis estadístico. \textbf{Verificó la detección de aleatoriedad de calidad criptográfica}, habilitando la evaluación rigurosa de generadores pseudoaleatorios (PRNG).}

\newcommand{\pixellabOpeningSecurity}{Desarrolló un kit de herramientas en Python para validación de PRNG. \textbf{Logró detección automatizada de fuentes de aleatoriedad débil mediante análisis de entropía y múltiples métricas estadísticas}.}

\newcommand{\pixellabOpeningSoftware}{Desarrolló un paquete Python para generación de imágenes a nivel de píxel y criptoanálisis.}

\newcommand{\pixellabOpeningApplied}{Implementó un kit de análisis basado en teoría de la información que demuestra la aplicación práctica de criptoanálisis estadístico y metodología de pruebas de aleatoriedad. \textbf{Validó métricas teóricas mediante estimación de $\pi$ por Monte Carlo y análisis de correlación en más de 50,000 casos de prueba}, conectando la teoría criptográfica con la validación empírica.}

\newcommand{\pixellabStatsCrypto}{Implementó métricas criptanalíticas integrales que incluyen cálculo de entropía de Shannon, pruebas de uniformidad chi-cuadrada y análisis de correlación multidireccional.}

\newcommand{\pixellabStatsSecurity}{Construyó un motor de análisis estadístico que mide la calidad de aleatoriedad mediante entropía, chi-cuadrada, correlación y métodos de Monte Carlo.}

\newcommand{\pixellabStatsSoftware}{Implementó manipulación a nivel de píxel \textbf{con soporte para asignación directa, generación recursiva y procesamiento por lotes}.}

\newcommand{\pixellabStatsApplied}{Desarrolló un marco de interpretación de métricas con herramientas de visualización para histogramas de entropía, mapas de calor de correlación y análisis de distribución. \textbf{Creó ejemplos educativos que demuestran conceptos de teoría de la información mediante generación interactiva de imágenes}, haciendo accesibles conceptos teóricos complejos.}

\newcommand{\pixellabArchitectureCrypto}{Diseñó una arquitectura modular con clases separadas ImageGenerator e ImageMetrics. \textbf{Logró un diseño de API limpio que permite la integración en marcos de pruebas de seguridad más amplios}.}

\newcommand{\pixellabArchitectureSecurity}{Creó una suite de pruebas integral que valida la correctitud de los algoritmos estadísticos contra casos de prueba conocidos.}

\newcommand{\pixellabArchitectureSoftware}{Construyó infraestructura de desarrollo completa que incluye pruebas automatizadas, herramientas CLI y scripts de ejemplo.}

\newcommand{\pixellabArchitectureApplied}{Publicó documentación exhaustiva que incluye referencia de API, ejemplos de uso y fundamentos teóricos sobre métricas criptanalíticas. \textbf{Creó más de 4 ejemplos educativos que demuestran pruebas de PRNG, arte algorítmico y análisis de esteganografía}, apoyando aplicaciones de investigación y educación.}

% =============================================================================
% Thesis project bullets
% =============================================================================
\newcommand{\thesisBulletCrypto}{Diseñó un esquema de comunicación cuántico-resistente de extremo a extremo que integra componentes de intercambio de claves y cifrado de datos, con fundamentos matemáticos rigurosos que respaldan tanto la correctitud funcional como las pruebas de seguridad bajo modelos de amenaza post-cuánticos.}

\newcommand{\thesisBulletSecurity}{Diseñó un esquema de comunicación cuántico-resistente de extremo a extremo que integra componentes de intercambio de claves y cifrado de datos para investigación de tesis de maestría, abordando las amenazas de seguridad emergentes de la computación cuántica.}

\newcommand{\thesisBulletSoftware}{Diseñó un sistema de comunicación modular con clara separación de responsabilidades entre las capas de intercambio de claves, cifrado y comunicación, enfatizando la arquitectura limpia y la reutilizabilidad de componentes para investigación de tesis de maestría.}

\newcommand{\thesisBulletApplied}{Diseñó un protocolo de comunicación cuántico-resistente que combina el diseño criptográfico teórico con consideraciones prácticas de implementación, incluyendo análisis formal de seguridad e implementación de prueba de concepto para tesis de maestría.}

% =============================================================================
% Educational Tools section content
% =============================================================================
\newcommand{\edutoolsPlatformDesc}{\textbf{Plataforma:} Sitio web público que muestra investigación criptográfica, herramientas interactivas.}

\newcommand{\edutoolsDemoIntro}{\textbf{Demostraciones Criptográficas Interactivas:}}

\newcommand{\edutoolsDemoList}{%
	\item RSA: Generación de llaves, flujos cifrado/descifrado, esquemas de relleno (PKCS\#1, OAEP).
	\item Funciones Hash: Visualización del efecto avalancha, estimación de colisiones, cómputo de hash.
	\item ECC: Aritmética de puntos, protocolo de intercambio de llaves ECDH, firmas digitales ECDSA.%
}

% =============================================================================
% Professional Skills (soft skills) section — items
% =============================================================================
\newcommand{\softskillOne}{\textbf{Comunicación Técnica:} Presentó más de 21 seminarios técnicos sobre sistemas criptográficos.}

\newcommand{\softskillTwo}{\textbf{Investigación y Educación:} Elaboró contenido educativo para participación académica y pública.}

\newcommand{\softskillThree}{\textbf{Escritura Técnica:} Desarrolló documentación integral para proyectos de código abierto.}

\newcommand{\softskillFour}{\textbf{Aprendizaje Autónomo:} Dominó de forma independiente la criptografía post-cuántica y los estándares NIST.}

\newcommand{\softskillFive}{\textbf{Resolución de Problemas:} Optimizó implementaciones criptográficas mediante optimización algorítmica y programación de bajo nivel.}

\newcommand{\softskillSix}{\textbf{Colaboración Interdisciplinaria:} Conectó la criptografía teórica con el desarrollo de software práctico, trabajando eficazmente en los dominios matemático, de seguridad e ingeniería.}

% =============================================================================
% Conference Presentations section — items
% =============================================================================
% Note: conference titles are proper event names in Spanish; body is unchanged.
\newcommand{\conferenceOne}{\textit{``Sistema de cifrado asimétrico resistente a ataques de computadoras cuánticas''} (Asymmetric Encryption System Resistant to Quantum Computer Attacks), \textbf{Congreso Internacional CAECH}, CAECH \& THINSCEN, Mexico City, Mexico, October 2024.}

\newcommand{\conferenceTwo}{\textit{``Sistema de cifrado asimétrico resistente a ataques de computadoras cuánticas''}, \textbf{Congreso Interpolitécnico de Investigación para Alumnos de Posgrado}, Instituto Politécnico Nacional, Mexico City, Mexico, June 2024.}

% =============================================================================
% Education section — entries
% =============================================================================
\newcommand{\educationOne}{%
\textbf{Maestría en Ciencias de la Computación} $|$ Centro de Innovación y Desarrollo Tecnológico en Cómputo, Instituto Politécnico Nacional, Ciudad de México $|$ Enero 2023 -- Julio 2025\\
\textit{Especialización: Criptografía y Sistemas Seguros}%
}

\newcommand{\educationTwo}{%
\textbf{Licenciatura en Física y Matemáticas} $|$ Escuela Superior de Física y Matemáticas, Instituto Politécnico Nacional, Ciudad de México $|$ Agosto 2017 -- Julio 2022\\
\textit{Especialización: Matemáticas Aplicadas}%
}
